\section{动态放置机制面临的关键挑战}
\label{sec:challenges}

\subsection{动态路由与专家负载变化的适应性}

路由局部性指同一样本 token 的 expert 选择关联性，网络局部性指样本所在节点或设备与目标 expert 的拓扑邻近程度；两者、expert popularity 和跨节点通信占比均随窗口变化。上一窗口的热点不一定延续，路由局部性也可能与全局专家负载趋势不同；调整过慢会滞后，逐步调整又可能暴露求解和状态更新成本。现有方法分别利用局部性、popularity 或历史工作负载预测，但缺少在同一比较条件下判断两类信号持续性的统一标准。

观测粒度也不同：路由局部性按样本或节点统计，网络局部性按样本设备与目标专家的拓扑关系统计，专家 popularity 按专家和窗口统计，窗口长度并不天然相同。

\subsection{动态调整收益与重配置成本的权衡}

判断是否值得调整，需要估计窗口内能减少多少关键路径等待，以及哪些成本无法被计算隐藏。样本放置的求解和索引、专家副本的权重与梯度同步、SYMI 的优化器状态位置，以及 Janus 的拉取和缓存，都可能改变端到端 step time。现有论文指标各异，缺口在于把暴露等待和调整成本放入同一观察窗口。

若策略还改变 expert capacity 或 token dropping，单步时间不能代表训练总成本：更低延迟可能伴随更多 token 丢弃，最终需要更多迭代达到同一质量。因此还应观察 per-step 暴露时间和 time-to-target-quality。

\subsection{通信局部性与专家负载的协同优化}

针对单一瓶颈的优化可能把等待转移到另一处：减少跨节点字节却未减轻热门专家计算，或复制后计算更均衡却新增副本同步。因此需要研究：系统能否识别主导的是网络路径、专家计算还是调整控制，并在保持现状、样本重排和专家复制之间进行有成本意识的选择？这连接了第 2 章的机制差异与下一章的个人方案。

